\subsection{Heterogeneity in delay valuation} 
\label{subsec:essay2_heterogeneity} 

While the overall MWTA provides a useful benchmark, substantial heterogeneity emerges across demographic, behavioral, and institutional dimensions. Table~\ref{tab:mwta_subgroups} reports subgroup-specific marginal willingness-to-accept (MWTA) for a one-month increase in vaccination delay, derived from multinomial logit models with vaccination delay and cash incentives entered linearly. Confidence intervals are cluster-robust at the respondent level and constructed using the delta method. 

For the overall sample, the MWTA is 44.8 RMB per month (95\% CI: 28.9--60.7), indicating that respondents require approximately this amount of compensation to tolerate an additional month of waiting. However, this average masks meaningful gradients across groups. 

\textit{Demographics.} Women exhibit a higher MWTA (55.7 RMB) than men (33.1 RMB), and urban residents display higher delay valuation (57.7 RMB) than rural residents (34.7 RMB). Education and age differences also emerge: lower-education respondents (51.2 RMB) and younger respondents (<55 years; 49.8 RMB) value timeliness more than higher-education respondents (28.1 RMB) and older respondents ($\geq$55 years; 30.8 RMB), although confidence intervals partially overlap. 

\textit{Health behavior.} Behavioral heterogeneity is particularly pronounced. Smokers require 55.0 RMB per month to accept delay, compared with 38.4 RMB among non-smokers. Differences across physical activity groups are substantial: individuals with low physical activity exhibit the highest MWTA (80.1 RMB; 95\% CI: 44.6--115.6), nearly three times that of the mid-activity group (27.6 RMB; 95\% CI: 5.1--50.1). 

\textit{Institutional trust.} The most striking contrast appears along institutional trust. Consistent with the categorization used in Essay 1, respondents are partitioned into \textit{High Institutional Trust} (Strongly/Somewhat trust) and \textit{Low Institutional Trust} (Neutral/Distrust) cohorts. Respondents in the low-trust group display an MWTA of 62.3 RMB (95\% CI: 37.1--87.5), more than double that of high-trust respondents (28.7 RMB; 95\% CI: 6.5--50.8). This pattern is consistent with the conceptual framework in which lower trust heightens perceived uncertainty regarding delayed protection, thereby increasing the behavioral burden of waiting. By grouping neutral responses with the distrust cohort, the analysis captures the behavior of the broader population lacking explicit institutional confidence, who bear significantly higher shadow prices for vaccine delay.

\textit{Medical condition and self-rated health.} Individuals reporting a medical condition exhibit higher MWTA (57.5 RMB) relative to those without (43.4 RMB), though the confidence interval is wider due to the smaller subgroup size. Differences by self-rated health are comparatively modest (44.8 RMB vs 44.9 RMB). 

Figure~\ref{fig:mwta_forest} provides a compact visualization of these subgroup differences. Point estimates and confidence intervals are plotted against a vertical reference line marking the overall MWTA, making clear which groups place above-average monetary value on timeliness. The figure highlights especially elevated delay costs among low-trust respondents and those with low physical activity, reinforcing the conclusion that the welfare burden of delay is unevenly distributed. 

Taken together, these results indicate that the behavioral cost of delay is far from uniform. Delay schedules that appear moderate at the aggregate level may impose disproportionate welfare burdens on specific subpopulations, particularly those with lower institutional trust or lower preventive engagement. These heterogeneity patterns motivate targeted prioritization or compensation strategies, explored further in the policy simulations below.
